%! Equations and Graphs in Both Directions — Unit 2, Lesson 2.3 — Student Packet Cover Sheet
\documentclass[10pt]{article}
\usepackage{atda-article}
\usepackage{atda-boxes}
\usepackage{ltablex}
\keepXColumns

\begin{document}

% Full-bleed royal banner
\begin{tikzpicture}[remember picture, overlay]
  \fill[royal] (current page.north west)
    rectangle ([yshift=-0.9in]current page.north east);
\end{tikzpicture}
\vspace{-0.8in}

\begin{center}
  {\color{white}\textbf{\LARGE Algebra, Trigonometry, and Data Analysis}}\\[4pt]
  {\color{mistmid}\large\textbf{Unit 2: Functions, Transformations, and Graph Analysis}}\\[6pt]
  {\color{white}\textbf{\large Lesson 2.3 \quad Equations and Graphs in Both Directions}}
\end{center}

\vspace{0.22in}
\namedateperiod
\vspace{0.18in}

\begin{learningtargetbox}
\small
\begin{enumerate}[leftmargin=1.5em, itemsep=5pt]
  \item \ldots{} \textbf{graph} a linear, quadratic, or exponential function from its equation by
        moving the parent function's \textbf{anchor point} and carrying a few points with it.
  \item \ldots{} \textbf{write the equation} of a function from its graph --- reading the vertex, the
        horizontal asymptote, or two points, depending on the family.
  \item \ldots{} find the stretch factor $a$ from one extra point when the image is not the same
        shape as the parent.
  \item \ldots{} use a graphing calculator to \textbf{verify} a transformation: type it correctly,
        set a window that actually shows the graph, and read the \textsc{table}.
  \item \ldots{} settle a disagreement between a picture and an equation by \textbf{substituting one
        input}, and explain why a matching screen is not proof.
\end{enumerate}
\end{learningtargetbox}

\vspace{0.14in}

\begin{tocbox}
\small
\renewcommand{\arraystretch}{1.5}
\begin{tabularx}{\linewidth}{@{} c l X r @{}}
\rowcolor{mist}
\textbf{\#} & \textbf{Component} & \textbf{Description} & \textbf{Score} \\
\midrule
1 & Warm-Up        & Two students type almost the same thing and get different graphs
                                                                    & \blank{1.2cm} \\
2 & Guided Notes   & Equation $\to$ graph, graph $\to$ equation, and what the TI-84 is actually
                     good for                                       & \blank{1.2cm} \\
3 & Group Activity & Building graphs, reading equations off them, and finding two typos
                                                                    & \blank{1.2cm} \\
4 & Exit Ticket    & One equation graphed, one equation written, and one screen settled by
                     substitution                                   & \blank{1.2cm} \\
5 & Homework       & The yearbook's profit curve, and a look ahead
                                                                    & \blank{1.2cm} \\
\midrule
  & \multicolumn{2}{r}{\textbf{Total}} & \blank{1.2cm} \\
\end{tabularx}
\end{tocbox}

\vspace{0.14in}

\begin{remindbox}
\small
This lesson covers \texttt{AFDA.AF.1d} (write the equation of a function \emph{given a graph},
using transformations of the parent) and \texttt{AFDA.AF.1f} (graph a function \emph{given its
equation}, using transformations --- and use technology to verify). Lesson 2.2 asked you to
\emph{describe} a move somebody showed you. Today you have to \emph{produce} the other half, in
both directions. One sentence runs the whole packet: \textbf{a picture can look right and be wrong,
but a substitution gives a number, and a number settles it.} The calculator draws exactly what you
typed --- never what you meant.
\end{remindbox}

\end{document}
